\section{Core module}
This section will focus on the objects in the core module. This module contains
the most vital objects in Snowman, including rules, rule revisions and sensors.

Simple objects such as generators, rule references and classtypes, which are
connected to rule objects, are not described in this section. For an overview of
all objects and functions in this module, refer to class diagram (ref) or
appendix (ref).

\subsection{Rules and rule revisions}
Rules in Snowman are split into two parts; the rule and the rule revision. As
threats change, or errors in the rules are discovered, a rule may become
obsolete and a new revision of the rule is released. The rulestring, message and
detection filter might change from one revision to the other, and are stored in
the rule revision object along with the revision number. A revision is always
connected to a rule object which contains data that is constant for all its
revisions, such as signature ID and rule set. The reason behind keeping rules
and revisions in two separate objects is to avoid writing the same constant rule
data for all revisions. Users of Snowman have the possibility to control how
many revisions to store per rule, and whether incoming revisions should be
active or inactive by default. This feature allows users to control
automatically downloaded rules before they are put into the detection
system. Users can also manually activate or deactivate rule revisions, and
Snowman will distribute only the newest active revision of a rule.

\subsection{Rule sets}
As previously stated, the categorization of rules in rule sets is a de facto way
of organizing rules. Rules in the community rule set\footnote{``The Community
Ruleset is a GPLv2 VRT certified ruleset that is distributed free of charge
without any VRT License restrictions'' - http://www.snort.org/snort-rules/} have
a ``ruleset'' option within the ``metadata'' option in the rulestring which
specifies the rule set. In other releases, rules are organized into
files where the file name denote their rule set. These two ways of specifying
the rule set is supported by Snowman.

\subsection{Sensors}

\subsection{Comments}
The core module contains a comment object which can store
comments from users who wish to document their operations, e.g. the user can
attach a comment when adding a filter to a rule to describe why the filter was
added and which effects it will have on the system. Snowman also generate
comments for important events, such as \ldots